\section{Experiments}\label{sec:experiments}

The primary setting is training a forward pass that jumps: spiking thresholds, INT8
quantization, argmax attention, staircase activations and hard MoE routing, with
MAX-SAT adding an integer-valued objective at up to $10^6$ variables. The benchmark set tests the same optimizer across several discontinuous operators. MNIST, ETTh1 and
quantized RL policy search bound the claim from either side: two objectives where
exact gradients exist, and one where the gradient is exactly zero.
At matched optimizer steps \sysname{} leads all six gradient-free baselines on
validation accuracy across all $36$ architecture-baseline comparisons; matching the
evaluation budget reverses argmax attention, reverses one MAX-SAT size, and levels
hard MoE. Primary comparisons use five seeds unless stated otherwise. Evaluation-matched
argmax and hard-MoE comparisons use three; theory-mode MAX-SAT uses two, and the
theory-mode ladder uses one. Legacy ablations state their own sample counts and
selection protocol. Error bars denote one sample standard deviation.

%% ========================================================================
\subsection{Experimental Setup}\label{sec:setup}
%% ========================================================================

\paragraph{Datasets.}
\textbf{MNIST} \citep{lecun1998gradient}: 60K/10K handwritten digits (28$\times$28 grayscale).
\textbf{ETTh1} \citep{zhou2021informer}: 17,420 hourly oil temperature readings (8640/2880/2880 train/val/test split), z-score normalized using train statistics.
The SNN and non-differentiable model experiments use MNIST with model-specific
architectures, except argmax attention (Fashion-MNIST) and hard MoE (MNIST
and Fashion-MNIST combined into 20 classes, with the latter labels shifted by 10).
For each run seed, the primary image benchmarks reserve a random 10\% of the official
training data for validation: 54K/6K/10K train/validation/test images for each
single dataset and 108K/12K/20K for hard MoE. Methods share the split within
each seed. Inputs are divided by 255, then normalized using $(0.1307,0.3081)$
for MNIST and $(0.2860,0.3530)$ for Fashion-MNIST (mean, standard deviation).
No image augmentation is used. The controlled weighting comparison uses the fixed
split specified in Section~\ref{sec:controlled-comparison}. Dataset details are
in Appendix~\ref{app:datasets}.

\paragraph{Models.}
MNIST uses a two-layer MLP with 102K trainable parameters (MNISTNet).
Memory scaling uses a VGG-11-style spiking network (SNNVGG11Small, ${\sim}$2.4M parameters) with leaky integrate-and-fire neurons at $T_{\mathrm{seq}}$ timesteps.
The SNN accuracy benchmark uses SpikingMNISTNet, a smaller spiking network with hard LIF thresholds ($T_{\mathrm{seq}}{=}15$ timesteps) on standard MNIST.
Exact architectures are in Appendix~\ref{app:architectures}.

The accuracy SNN maps $784\to128\to10$ through two hard LIF layers.
At each of 15 steps, $\widetilde v_t=0.95v_{t-1}+I_t$,
$z_t=\mathbb{1}\{\widetilde v_t\geq1\}$, and
$v_t=\widetilde v_t(1-z_t)$, starting at zero. Its first-layer input current
is constant over time and the output is the sum of output spikes, used as
cross-entropy logits without division by 15. INT8 uses widths
$784\to128\to128\to10$ with ReLU and quantizes the middle layer's
weights and biases to $0.01\operatorname{clip}(\operatorname{round}(w/0.01),-128,127)$.
Staircase uses the same widths with two five-level staircase activations.
Argmax attention uses a 128-dimensional representation and eight slots.
Hard MoE uses a 128-dimensional representation, four experts with hard
argmax gating, and a 20-class output. The supervised classification objective
is mean cross-entropy on the current minibatch.

\paragraph{Baselines.}
Gradient-free:
\textbf{CMA-ES} \citep{hansen2016cma}, covariance matrix adaptation;
\textbf{OpenAI-ES} \citep{salimans2017evolution}, isotropic perturbations with
rank shaping;
\textbf{EGGROLL} \citep{sarkar2025eggroll}, low-rank evolution strategies;
\textbf{MeZO} \citep{malladi2023mezo}, in-place zeroth-order SGD with
seed-regenerated perturbations;
\textbf{SPSA} \citep{spall1992multivariate};
and \textbf{subspace random search}, which samples a random direction in the same
subspace and accepts on improvement. It shares \sysname{}'s representation;
its candidate set and weighting also differ. Section~\ref{sec:controlled-comparison}
holds the candidate law fixed across weighting rules while allowing their
natural displacement lengths to differ.
Gradient-based references where applicable: \textbf{Adam} \citep{kingma2015adam}
or SGD; \textbf{surrogate-gradient} BPTT \citep{neftci2019surrogate} on SNN tasks;
and \textbf{STE}/quantization-aware training \citep{bengio2013ste} on the
quantized tasks. These are reference implementations with backward approximations, not
upper bounds on achievable performance. The blocked-gradient controls use
neither a surrogate derivative nor STE.

\paragraph{Protocol.}
For the primary supervised results, every method trains on a training split, selects its
checkpoint on a validation split, and is scored once on test with that
checkpoint; the alternative, maximum test accuracy over epochs, is test-set
selection. On MNIST the two conventions differ by roughly $0.7$~pp for \sysname{} and
not uniformly across methods, so the difference does not cancel in a comparison.
Determinism, the pinned flags and the cross-device spread are in
Appendix~\ref{app:reproducibility}.

\paragraph{Hyperparameters.}
Every primary PolyStep architecture comparison uses the softmax solver (Eq.~\ref{eq:softmax-rule}), with
$K{=}1$ probe per vertex on the simplex, at $d_p{=}8$ in subspace mode and $d_p{=}2$ in
full-space mode. The column marginal is measured slack in every subspace run and
inconclusive on full-space MAX-SAT, the one primary setting where it could bind, where
the two solvers agree on the answer anyway (Appendix~\ref{app:ot-highparticle}). Every
method in these architecture comparisons, \sysname{} included, receives a two-stage grid of nine configurations per stage
at an equal per-configuration budget, selected on validation, the second stage
recentering each axis on the first stage's selection so that a boundary selection
extends the range instead of being accepted as tuned. Baselines start from reference
defaults with the probe scale converted to a per-coordinate quantity, since \sysname{}'s
radius is the norm of a displacement while a baseline's $\sigma$ is per coordinate;
using the same numerical value without this conversion would change the
expected perturbation norm by a factor of order $\sqrt d$. Where
\sysname{} anneals, the baselines hold the schedule's final value, which favors us.
Sweep mechanics, the selected configurations and the transfer rule are in
Appendix~\ref{app:hyperparameters}.

Within an update, all candidate losses use the same minibatch. The image
experiments use batches of 512, including the smaller final batch of an epoch.
A candidate evaluation is one parameter vector scored on that batch, regardless
of whether several candidates share a GPU call. Candidate counts therefore
differ from processed-example counts. Validation and final test evaluation are
separate from the reported training-candidate budgets. The selected settings,
epoch counts, and reuse intervals are in Table~\ref{tab:hyperparameters};
no test score is used to retune a primary configuration.

\paragraph{Compute and statistical testing.}
Every run is on one NVIDIA RTX 5090 (32\,GB) with software versions recorded in each result file, over
seeds $\{42, 123, 456, 789, 1337\}$. Pairwise comparisons use an exact two-sided
sign-flip permutation test over all $2^5$ seed assignments, whose smallest attainable
value is $2/2^5 = 0.0625$; that floor, and the absence of a multiplicity correction
across benchmarks, are in Appendix~\ref{app:stats}.

%% ========================================================================
\subsection{The Matching Axis}\label{sec:matching-axis}
%% ========================================================================

Optimizer steps, candidate evaluations and wall-clock rank the methods differently, so
every table names its axis. \emph{Optimizer steps} is the primary axis, a sequential
dependency no batching removes; one generation of a population method is one step, with population $32$ for the population baselines in this comparison. Table~\ref{tab:step-sweep} is that
comparison, across all $36$ pairings. The per-architecture tables instead give each baseline the
reported matching axis: the population methods run evaluation-matched on the SNN, MNIST, INT8 and
staircase, step-matched on argmax and hard MoE, and the finite-difference methods run
step-matched throughout, since at two evaluations per sequential step MNIST's budget
would demand $5{,}094{,}360$ of them. Each row carries the axis it was budgeted on.

Matching \emph{evaluations} unties the population and shrinks \sysname{}'s step count as
$1/d_{\mathrm{sub}}$: at MNIST's evaluation-matched budget OpenAI-ES takes \mnistStepRatio{} as
many optimizer steps. Section~\ref{sec:other-axes} gives that axis in full, including the
two architectures and the $5{,}000$-variable MAX-SAT size where it changes the answer.
Wall-clock measures practical training cost and depends on batching and hardware.
At population $32$ the measured throughput is about $90\times$ lower than with a
tuned population; Appendix~\ref{app:wallclock-matched} gives the timings, the
per-result bookkeeping that enforces the axis, and the one representation asymmetry,
EGGROLL's factored subspace. Otherwise every gradient-free method in a table shares the
subspace object, the minibatch stream, and a tuning sweep of the same shape.

Two conventions carry through. Table~\ref{tab:step-sweep} reports \emph{validation}
accuracy, since a trajectory truncated at another method's step count has no test score,
while every per-architecture table reports \emph{test} accuracy at the
validation-selected checkpoint; test runs $0.0$ to $0.9$~pp above validation here, so the
two differ in level and not in ordering. And where an evaluation-axis comparison ran three
seeds, its \sysname{} figure is the same run read on those three, not the five-seed figure.

%% ========================================================================
\subsection{The Benchmark Set}\label{sec:gallery-overview}
%% ========================================================================

The five hard architectures place discontinuities in different operators
(Table~\ref{tab:gallery-overview}). A hard operator can have zero derivative
between thresholds while the complete network still varies through other paths.
Surrogate-gradient and quantization-aware training provide alternative backward
rules; they need not change the hard forward model. MAX-SAT supplies a fully
discrete objective through its clause count.

A threshold crossing need not change the batch loss. The measured nonzero jumps
in Table~\ref{tab:gallery-overview} establish examples, rather than a uniform
lower bound across all walls. In the INT8 diagnostic, the batch jump lies
$13$--$300\times$ above the simple $J=j/n_B$ prediction over a $64\times$ range
of batch size. Appendix~\ref{app:discontinuity-sets} derives the local threshold
geometry, and Appendix~\ref{app:measured-jump} gives the measurement protocol.

\begin{table}[ht]
\centering
\caption{Non-differentiable benchmark gallery: local threshold geometry and observed nonzero loss jumps (zero-fraction of solved wall crossings and median $J$ shown; protocol in Appendix~\ref{app:measured-jump}). These measurements do not certify a uniform jump bound on every threshold sheet.}
\label{tab:gallery-overview}
\small
\begin{tabular}{lllccc}
\toprule
Benchmark & Discontinuity set $\mathcal{D}$ & Anchor & zero frac. & median $J$ & Section \\
\midrule
SNN (hard LIF spikes)            & spike-threshold surfaces & App.~\ref{app:discontinuity-sets}(a) & $0.69$ & $4.1{\times}10^{-7}$ & \ref{sec:nondiff} \\
INT8 quantization                & quantization-grid walls     & App.~\ref{app:discontinuity-sets}(b) & $0.11$ & $1.6{\times}10^{-5}$ & \ref{sec:nondiff} \\
Hard MoE (argmax routing)        & Voronoi walls               & App.~\ref{app:discontinuity-sets}(c) & $0.00$ & $1.2{\times}10^{-3}$ & \ref{sec:nondiff} \\
Staircase activations            & integer level sets          & App.~\ref{app:discontinuity-sets}(d) & $0.00$ & $2.9{\times}10^{-6}$ & \ref{sec:nondiff} \\
Argmax attention                 & score-tie surfaces       & App.~\ref{app:discontinuity-sets}(e) & $0.00$ & $1.5{\times}10^{-5}$ & \ref{sec:nondiff} \\
MAX-SAT 1M variables             & piecewise-constant          & Cor.~\ref{cor:piecewise-constant} & -- & -- & \ref{sec:maxsat} \\
\bottomrule
\end{tabular}
\end{table}

%% ========================================================================
\subsection{Non-Differentiable Model Training}\label{sec:nondiff}

\begin{figure}[t]
\centering
\includegraphics[width=\linewidth,keepaspectratio]{primary_summary}
\caption{Hard-LIF SNN on MNIST (five seeds, mean $\pm$ sample std).
(a)~Test accuracy at the validation-selected checkpoint; dots show individual
seeds and labels give evaluation budgets. Adam$^\dagger$ trains a surrogate
relaxation. (b,c)~Logged validation accuracy for all seven gradient-free
methods against optimizer updates and candidate evaluations. Curves linearly
interpolate logs within the five seeds' shared support. The dotted line marks
$3{,}180$ updates. Population methods receive $15.54$M evaluations;
MeZO, SPSA, and random search use the smaller budgets shown in (a).}
\label{fig:primary-summary}
\end{figure}
%% ========================================================================

Hard spike thresholds make SNNs a direct test of optimization from scalar
losses. A forward-only method can also train through a hardware interface
that supplies losses but no derivatives. We use SpikingMNISTNet with hard
LIF thresholds and $T_{\mathrm{seq}}=15$ timesteps on MNIST.

% SNN comparison table
\begin{table}[t]
\caption{SNN test accuracy (\%, MNIST, hard LIF, 5 seeds, mean $\pm$ std at the validation-selected checkpoint); best gradient-free result in bold. Population baselines are matched on \emph{evaluations}, so their step counts follow their populations; the finite-difference baselines are budgeted on steps. The two axes agree on this architecture. Adam$^\dagger$ trains a surrogate relaxation of the same architecture.}
\label{tab:snn-accuracy}
\centering
\small
\input{tables/snn_comparison}
\end{table}

\sysname{} leads the best tuned gradient-free baseline on every seed
(Table~\ref{tab:snn-accuracy}, Figure~\ref{fig:primary-summary}), at the design floor
$p = 0.0625$ of the exact sign-flip test, and does so on all five seeds of MNIST, INT8,
staircase, argmax attention and hard MoE as well.

Figure~\ref{fig:primary-summary} shows the selected test scores alongside
progress per update and per evaluation. \sysname{} reaches its reported accuracy in $3{,}180$ steps,
while the evaluation-matched population baselines take many more updates.
OpenAI-ES and CMA-ES make faster early progress per evaluation. The final
accuracy advantage therefore does not imply uniform superiority along the
learning curve or lower evaluation cost. MeZO and SPSA improve under their
much smaller evaluation budgets, but finish with lower selected test accuracy.

\paragraph{Surrogate-gradient training.}
The tuned Adam-surrogate baseline uses BPTT with a smooth spike
approximation~\citep{neftci2019surrogate}. It has lower mean test accuracy and
larger seed variation in this experiment. A separate slope sweep spans
$20$~pp across $\alpha\in[0.5,10]$, but it does not isolate the cause of that
variation. SLTT~\citep{meng2023towards} and OTTT~\citep{xiao2022online}
report accuracies above $99\%$ with different surrogate-gradient training
recipes. Our experiment establishes performance for the stated architecture
and tuning protocol; it does not close that gap. The selected \sysname{}
configuration uses fixed $\varepsilon=0.5$, $r_s=2.0$, and rotations biased
toward the previous update direction. It lies outside the convergence
assumptions in Section~\ref{sec:theory}.

\paragraph{The four remaining architectures.}
The same discontinuity sits in a quantizer bin, a floor level, a routing wall and a
score tie. Each is deployed for its own reason, integer inference on edge hardware,
integer-valued features, sparse expert routing and hard attention, and together they
test whether the result depends on one operator rather than on the jump.
\sysname{} trains all four (Appendix~\ref{app:gallery-tables}), and where a smooth
relaxation exists the surrogate-gradient Adam$^\dagger$ setting exceeds it by
$0.8$--$3.2$~pp, on INT8, argmax and staircase, the SNN being the exception. Binary and
ternary weight networks are not measured under this protocol.

\paragraph{Quantized head, pretrained backbone.}
This experiment isolates the quantizer from everything else: freeze a pretrained backbone,
train a small head, and switch only whether that head is quantized. We freeze GPT-2 124M
and train its $1{,}538$-parameter SST-2 head, with every gradient-free method tuned by
the same nine-configuration validation sweep at a matched $500$K-evaluation budget
(Appendix~\ref{app:gpt2-headquant}). The ordering reverses with the quantizer
(Table~\ref{tab:gpt2-headquant}).

The head receives cached final-token features from the frozen backbone.
Of 5000 SST-2 training examples, a seed-0 permutation assigns 4500 to
updates and 500 to validation; the official 872-example validation set is
used as the held-out test set because the official test labels are unavailable.
The nine-configuration sweep uses 100K candidate evaluations per setting;
the selected configuration then receives 500K evaluations on each of five
seeds. Validation selection considers both the mean iterate and the best
sampled candidate. INT8 quantizes weight and bias at scale 0.01, binary
quantizes only weight by its sign, and the smooth control quantizes neither.
The selected head is tested once. The separate legacy all-parameter and
head-only GPT-2 experiments in the appendix use a different selection rule
and are not pooled with this comparison.

On the smooth control, MeZO and EGGROLL lead PolyStep. A MeZO update costs
$2$ evaluations against PolyStep's approximately $2{,}300$, permitting many more
updates at the same evaluation budget. Quantization changes the ordering, but this
comparison does not isolate estimator quality from update count or tuning.
MeZO's INT8 sweep selects a radius spanning about one hundred quantization bins;
such finite-radius differences can still estimate a smoothed gradient.
Adam's gradient through the INT8 head is zero. In the binary head, the two bias
parameters remain differentiable. The flat-pair diagnostic of
Section~\ref{sec:gallery-overview} uses this head.

\begin{table}[t]
\centering
\caption{GPT-2 SST-2 head fine-tuning through a quantizer (frozen 124M backbone, 1,538-parameter head; test accuracy \%, 5 seeds, mean $\pm$ std at the validation-selected checkpoint). Gradient-free methods are matched at $500$K evaluations and tuned by the same nine-configuration validation sweep; best gradient-free result per column in bold. Adam's gradient through the INT8 head is exactly zero; the binary bias (2 of 1,538 parameters) stays differentiable, a partial control.}
\label{tab:gpt2-headquant}
\small
\begin{tabular}{lccc}
\toprule
Method & INT8 & Binary & Smooth control \\
\midrule
\sysname{} & $67.1 \pm 6.9$ & $\mathbf{72.7 \pm 3.3}$ & $51.2 \pm 1.4$ \\
MeZO & $\mathbf{67.3 \pm 10.9}$ & $52.0 \pm 3.5$ & $62.6 \pm 2.5$ \\
EGGROLL & $65.8 \pm 2.5$ & $71.7 \pm 4.0$ & $\mathbf{63.7 \pm 1.0}$ \\
\midrule
Adam & $49.8 \pm 1.1$ & $50.1 \pm 1.1$ & $81.6 \pm 0.6$ \\
\bottomrule
\end{tabular}
\end{table}

\begin{table}[t]
\caption{Every method on every architecture at \emph{matched optimizer steps}: each baseline read at the step count \sysname{} spends ($3{,}180$; $6{,}330$ on hard MoE). Validation accuracy (\%), five-seed means (a truncated trajectory carries no test score); standard deviations in Table~\ref{tab:snn-accuracy} and Appendix~\ref{app:gallery-tables}. MNIST is the one column where \sysname{} reuses a transport plan and so evaluates on every third step (Section~\ref{sec:ablation}).}
\label{tab:step-sweep}
\centering
\small
\input{tables/step_sweep}
\end{table}

\paragraph{The step axis.}
Table~\ref{tab:step-sweep} places all six architectures and all six gradient-free
baselines on the sequential-step axis. A run budgeted on steps keeps its recorded
selection; a run budgeted on evaluations is read off its trajectory at the best
validation point on or before that step count. The table therefore answers progress per
sequential step, not progress per evaluation with the population untied. Sparse-expert systems
\citep{fedus2022switch,lepikhin2021gshard,zhou2022expertchoice} use different
routing and training recipes. The present benchmark does not establish a
state-of-the-art comparison with those systems.

\subsubsection{The Evaluation Axis}\label{sec:other-axes}

Untying the budget leaves four of the six architectures unchanged and reverses the
other two.

\emph{Four are decided before any budget binds.} On the SNN, MNIST, INT8 and staircase
the baselines stop improving long before their evaluations run out, peaking between
$0.2\%$ and $25\%$ of budget and declining or flattening after; a larger population
reaches the same ceiling sooner and no higher. Against the closest baseline in each,
validation margins run $2.5$ to $14.0$~pp at matched evaluations and $7.0$ to
$15.3$~pp at matched steps (Table~\ref{tab:axes}).

\emph{On argmax attention \sysname{} loses.} Given \sysname{}'s $42.01$M evaluations
and a population able to spend them ($2{,}048$, learning rate retuned on validation),
OpenAI-ES takes $20{,}514$ steps and reaches $86.95\%$ test over seeds $42/123/456$
against \sysname{}'s $86.41\%$ read on the same three, leading on every paired seed in a
fifth of the wall-clock. The matched-step margin, $6.6$~pp on test and $7.1$~pp on the
validation figure Table~\ref{tab:step-sweep} prints, is carried by the step axis alone.

\emph{On hard MoE the two methods are level.} At the full $100.55$M evaluations the
same configuration reaches $90.52\pm0.11\%$ against \sysname{}'s
$90.49\pm0.13\%$ on the same three seeds. MoE is where the
step cap supplied the largest advantage, $496\times$ the evaluations, and where
removing the cap removes the result.

\emph{The step-deficit diagnosis fails a direct test.} If the argmax result were a step
deficit, more steps should recover it. \sysname{}'s counterpart to the population size
is the subspace rank, which trades search dimension against step count at a fixed budget,
and accuracy is monotone in rank: halving the subspace yields $1.9\times$ the steps at a
$1.0$~pp cost and quartering it yields $3.7\times$ at $2.4$~pp
(Appendix~\ref{app:wallclock-matched}). The reported configuration is the best of the
three. This rules out those two rank reductions as remedies; it does not
exclude other tuning choices.

\paragraph{Memory scaling.}\label{sec:snn-memory}
Forward-only evaluation avoids retaining the activation graph used by
ordinary BPTT. On SNNVGG11Small at batch size $4$, the measured forward-only
peak rises from $31.8$ to $51.6$\,MiB as $T_{\mathrm{seq}}$ increases from
$25$ to $400$; forward plus BPTT rises from $132.3$ to $1{,}538.2$\,MiB
(Figure~\ref{fig:scalability-plate}b). The ratio is $29.8$ at the longest
horizon. These measurements isolate evaluation memory. They exclude a
complete \sysname{} step's candidate and assignment buffers and do not
establish an asymptotic memory law.

%% ========================================================================
\subsection{Discrete Optimization at Scale: MAX-SAT}\label{sec:maxsat}
%% ========================================================================


MAX-SAT is the limiting case: the objective is integer-valued, so no smoothing of any
one operator recovers a derivative, and the instance size can be pushed until a
population no longer covers the search space. We run random 3-SAT at the critical ratio
$\alpha{=}4.27$ from $100$ to $10^6$ variables, on $\sigma(\mathbf{x})$ with hard
rounding $\lfloor\cdot\rceil$. At $10^6$ variables the instance has $4.27$M clauses,
our largest non-differentiable run.

\begin{table}[t]
\centering
\caption{Random MAX-3-SAT, $\alpha=4.27$, final clauses satisfied
(\%). One fixed instance per size; general-purpose methods use five
optimizer seeds (mean $\pm$ sample std), and domain-specialized
references use one run.
\sysname{} and OpenAI-ES use $1{,}000$ updates below $10^6$ variables
and $500$ at $10^6$. CMA-ES receives OpenAI-ES's evaluation budget
($50{,}000$ below $10^6$), so its update count varies; it uses diagonal
covariance from $1{,}000$ variables. Matching \sysname{}'s evaluations
at $5{,}000$ variables reverses the ordering
(Section~\ref{sec:matching-axis}). Best per block is bold.}
\label{tab:maxsat-scaling}
\small
\begin{tabular}{lccc@{\hspace{0.4em}}|@{\hspace{0.4em}}ccc}
\toprule
 & \multicolumn{3}{c}{\textbf{General-purpose GFO}} & \multicolumn{3}{c}{\textbf{Domain-specialized}} \\
\cmidrule(lr){2-4} \cmidrule(lr){5-7}
\textbf{Variables} & \textbf{\sysname{}} & \textbf{CMA-ES} & \textbf{OpenAI-ES} & \textbf{probSAT}$^{\ddagger}$ & \textbf{SLS} & \textbf{RC2}$^\dagger$ \\
\midrule
100         & $98.0 \pm 0.7$ & $99.0 \pm 0.4$ & $\mathbf{99.4 \pm 0.1}$ & $\mathbf{99.8}$ & $\mathbf{99.8}$ & $\mathbf{99.8}$ \\
500         & $98.1 \pm 0.2$ & $\mathbf{98.4 \pm 0.3}$ & $98.2 \pm 0.2$ & $\mathbf{100.0}$ & $99.9$ & timeout \\
1{,}000     & $\mathbf{98.1 \pm 0.2}$ & $97.3 \pm 0.1$ & $96.9 \pm 0.2$ & $\mathbf{100.0}$ & $99.9$ & timeout \\
5{,}000     & $\mathbf{98.2 \pm 0.1}$ & $95.2 \pm 0.2$ & $92.9 \pm 0.2$ & $\mathbf{99.9}$ & $99.6$ & timeout \\
20{,}000    & $\mathbf{98.2 \pm 0.0}$ & $92.9 \pm 0.2$ & $90.5 \pm 0.1$ & $\mathbf{99.8}$ & $99.1$ & timeout \\
100{,}000   & $\mathbf{98.1 \pm 0.01}$ & $90.7 \pm 0.1$ & $88.9 \pm 0.0$ & $\mathbf{99.6}$ & $97.4$ & timeout \\
$10^6$$^{\S}$ & $\mathbf{92.6 \pm 0.02}$ & -- & $87.8 \pm 0.0$ & $\mathbf{98.9}$ & $89.4$ & timeout \\
\bottomrule
\end{tabular}
\vspace{0.3em}

{\raggedright\scriptsize $^\dagger$RC2 reaches its time limit for
$n\ge500$. $^{\ddagger}$probSAT uses ten restarts, up to $10^7$ flips
($5\times10^6$ at $10^6$ variables), with a wall-time allowance tied to
\sysname{}; it finishes in $13$--$89$ seconds.
$^{\S}$The million-variable results use a separate configuration, with incremental evaluation for \sysname{}. CMA-ES was not
run at that size.\par}
\end{table}

This is a scaling stress test of forward-only optimization, not a competitor to
dedicated SAT solvers: probSAT~\citep{balint2013probsat} leads at every scale and the
table marks that ceiling. A uniformly random assignment satisfies $7/8=87.5\%$ of clauses in
expectation, the reference value \citep{johnson1974approximation}, and staying clear of it
as the variable count scales is what a vanishing per-step cost difference makes hard.

Two regimes separate at $1{,}000$ variables (Table~\ref{tab:maxsat-scaling}). Below it
the evolution strategies lead by a few clauses, about six of $427$ at the widest such
margin. From $1{,}000$ to $100{,}000$ variables, \sysname{} stays near $98.1\%$
while both evolution strategies approach the random-assignment reference.
The separate million-variable configuration reaches $92.6\%$.
The budget matters: at $5{,}000$ variables, OpenAI-ES reaches $98.96\%$
against \sysname{}'s $98.15\%$ when given \sysname{}'s evaluation
budget (Figure~\ref{fig:maxsat-scaling-full}d). CMA-ES uses diagonal covariance from $1{,}000$ variables and
is not evaluated at $10^6$; limited-memory variants
\citep{loshchilov2017lmcma,akimoto2020ddcma} are outside this comparison.
The spread measures optimizer variation on a fixed instance, not
variation across new instances; it does not identify the mechanism
behind the observed accuracy.

The million-variable solve averages $1{,}216$ seconds on one RTX~5090
with $2{,}310.5$\,MiB peak GPU memory. Its incremental evaluator and
$500$-update budget differ from the $100{,}000$-variable runs, which
average $9{,}236$ seconds over $1{,}000$ updates with the full evaluator
(Figure~\ref{fig:maxsat-scaling-full}, Appendix~\ref{app:scalability}).
These timings describe the specified implementations and budgets.

%% ========================================================================
\subsection{Theory Mode}\label{sec:theory-mode}
%% ========================================================================

None of the tuned configurations satisfies the schedule condition,
Assumption~\ref{asm:main}(v) (Table~\ref{tab:hypothesis-coverage}). To measure what
that costs, we run \emph{theory mode}, with the following settings:
flat $\varepsilon$, step radius $r_{s,t} = r_0 (t+1)^{-0.6}$ so that
$\sum r_{s,t} = \infty$ and $\sum r_{s,t}^2 < \infty$, independently sampled
rotations, mollifier jitter meeting condition~(iv)'s $W^{2,1}$ requirement
(Lemma~\ref{lem:smooth-kernel}), the orthoplex
frame, HybridSubspace, and none of the acceleration heuristics of
Appendix~\ref{app:turbo-mode}. Every other knob is held at the tuned value, so the
comparison measures these configuration changes at the retained settings. It does
not isolate a necessary cost of satisfying the theorem. The retained sparse
bases and epoch-shuffled minibatches still fail orthonormality and conditional
unbiasedness, respectively (Section~\ref{sec:hypothesis-coverage}).
The orthoplex is the one restriction the theorem does not require (condition~(i)
admits the simplex the tuned runs use), and it makes theory mode spend more
evaluations per step than the row it is compared against (Section~\ref{sec:ablation}).

\begin{table}[t]
\centering
\caption{Theory mode versus the tuned configuration. Frame, rotation, jitter, and schedule settings are prescribed; projection, sampling, analytic assumptions, and padding remain qualifications (Section~\ref{sec:hypothesis-coverage}). SNN and MNIST use 5 seeds and report test accuracy at the validation-selected checkpoint. Theory-mode MAX-SAT uses 2 seeds (42 and 123) and reports final clause satisfaction; its tuned comparison uses 5 seeds. The gap is descriptive, not a causal decomposition of the theorem assumptions.}
\label{tab:theory-mode}
\small
\begin{tabular}{@{}lccc@{}}
\toprule
Benchmark & Theory mode & Tuned & Gap \\
\midrule
SNN (hard LIF) & $54.7 \pm 16.0$\% & \snnBestAcc{} & $38.3$~pp \\
MNIST          & $91.7 \pm 0.3$\% & \mnistBestAcc{} & $5.1$~pp \\
MAX-SAT (100K) & $96.9 \pm 0.1$\% & \maxsatHundredKAcc{} & $1.2$~pp \\
\bottomrule
\end{tabular}
\end{table}

\paragraph{Individual configuration changes.}
Theory mode changes five things at once, so Table~\ref{tab:theory-ladder} switches the
four listed settings one at a time in both directions. In this single-seed comparison, the largest change is associated with
condition~(v)'s decaying step radius. The other changes are smaller. The fifth change, disabling the acceleration heuristics, has no ladder row of
its own, since probe jitter and plan reuse are mutually exclusive
(Section~\ref{sec:ablation}), so it is folded into the jitter rows and the residuals.
The effects depend on the starting configuration and do not sum to the
full difference. The ladder's seed-42 theory-mode result lies within the
five-seed spread of Table~\ref{tab:theory-mode}.

\begin{table}[ht]
\centering
\caption{Effects of individual configuration changes toward Assumption~\ref{asm:main} (SNN, seed 42, test accuracy at the validation-selected checkpoint). Upper block starts from the tuned configuration and adds one hypothesis; lower block starts from theory mode and removes one. Each $\Delta$ is against the first row of its own block, not against the row above.}
\label{tab:theory-ladder}
\small
\input{tables/theory_ladder}
\vspace{0.3em}

{\raggedright\scriptsize Conditions~(iii), (vi), (vii) and~(viii) are not flippable:
(iii) is on in both, (vi) and~(viii) are properties of the trajectory, not
settings, and (vii)'s transverse half is a loss-regularity assumption
(Table~\ref{tab:hypothesis-coverage}).\par}
\end{table}

\paragraph{Finite-horizon behavior.}
Over $3{,}180$ steps, the schedule $r_{s,t}=r_0(t+1)^{-0.6}$ reduces the
step multiplier by a factor of $126$. The displacement upper bound shrinks
with it; the lower envelope in Proposition~\ref{prop:fragility} additionally
requires concentration. In a seed-42 check, multiplying $r_0$ by
$2$, $4$, $10$, and $100$ gives $56.2\%$, $47.6\%$, $9.8\%$, and $20.8\%$,
respectively, compared with $65.4\%$ at the ladder's original $r_0$.
These settings do not recover the tuned result. They do not rule out other
schedule choices or isolate decay as the sole cause of the gap.
The horizon-dependent constant step multiplier of
Corollary~\ref{cor:finite-horizon} is a distinct theoretical schedule.

\subsection{MNIST}\label{sec:vision}
%% ========================================================================

MNIST is a smooth benchmark, a setting favorable to zeroth-order methods, and serves here as a sanity check.

% MNIST comparison table
\begin{table}[t]
\caption{MNIST test accuracy (\%, 5 seeds, mean $\pm$ std at the validation-selected checkpoint), with both cost axes and the axis each row was budgeted on; best gradient-free result in bold.}
\label{tab:mnist-accuracy}
\centering
\small
\input{tables/mnist_comparison}
\end{table}

\sysname{} trails Adam by $0.9$~pp (Table~\ref{tab:mnist-accuracy}). The configuration is a cosine temperature schedule at
subspace rank~8 with no momentum and transport-plan reuse every three steps; momentum
costs $0.9$~pp here, on an objective smooth enough that the probes already produce
informative cost differences. MNIST also appears elsewhere under other configurations,
each isolating one factor; only the primary $30$-epoch subspace result is a claim.
Every gradient-free method plateaus below Adam, which is the information asymmetry
between a $d$-dimensional gradient and a scalar.

%% ========================================================================
\subsection{Reinforcement-Learning Policy Search}\label{sec:rl-policy-search}
%% ========================================================================

Forward-only policy search uses scalar episodic returns and can train
policies with hard quantization. On CartPole-v1 and Acrobot-v1,
\sysname{} reaches $0.995$--$1.000$ normalized final return across Float32,
INT8, and binary hidden activations. It leads OpenAI-ES under a shared
$200$-update budget, with unequal rollout costs. The normalized metric
clips at one, so it can conceal return differences near the reference.
PPO and DQN provide differentiable controls and additional diagnostics
with hidden and output quantization, which blocks their gradients without
a straight-through estimator. The different quantization placement limits
comparisons between these controls and the forward-only methods.
Appendix~\ref{app:rl-policy-search} gives the curves, costs, and protocol;
Proposition~\ref{prop:rl-policy-search} controls the effect of rollout
noise on assignment weights.

Both forward-only policies have one hidden layer of width 16. Float32 uses
ReLU; INT8 rounds the first hidden activation to a grid with scale
$s=\max|x|/127$, and binary uses its sign. The final policy is evaluated
without validation-based checkpoint selection. PPO/DQN instead quantize all
hidden activations and their outputs without STE, use additional
algorithm-specific networks, and receive 1M training transitions on CartPole
and 0.5M on Acrobot. Curves retain raw returns; the table clips normalized
scores using $(R-22)/478$ and $(R+500)/420$, respectively. The reported
rollout counters differ: PolyStep counts realized transitions, ES includes
allocated but unused horizon slots, and PPO/DQN include final evaluation
allowances. They are not a common sample-efficiency measure.

%% ========================================================================
\subsection{LSTM Time-Series Forecasting}\label{sec:lstm-timeseries}
%% ========================================================================

ETTh1 provides a smooth regression control with exact gradients. We train a
$23$K-parameter LSTM with a $96$-step lookback and forecast horizon
(Table~\ref{tab:lstm-timeseries}, Appendix~\ref{app:lstm-details}).
PolyStep matches Adam on validation error. On test, PolyStep, Adam, and
EGGROLL are close relative to their seed variation, with EGGROLL lowest on
the point estimate. Every learned method trails persistence. A level shift
between train and test may contribute, but this experiment does not isolate
its cause. PolyStep uses a hand-selected configuration and the baselines use
reference defaults; the two-stage architecture sweep does not cover this task.
The chronological split uses observations 0--8639 for training, 8640--11519
for validation, and 11520--14399 for testing. Normalization uses training
statistics only. The final 3020 observations are outside this experiment.

\subsection{Observed Training Instability}\label{sec:instability}

Some runs lose accuracy late in training, and the schedule ablation shows
larger drops under the tested temperature-decay settings. The concentration
bound in Proposition~\ref{prop:fragility} controls a single displacement;
the quadratic drift calculation has additional hypotheses. Neither result
predicts an accuracy drop or proves instability for the trained networks.
Likewise, smoothing blow-up is a local seminorm bound, not a lower bound
on every realized update near a threshold. The fixed-temperature,
decaying-step schedule is the variant analyzed by
Theorem~\ref{thm:piecewise-smooth}. The empirical schedule comparisons in
Appendix~\ref{app:schedule-fragility} describe the tested configurations.

\subsection{Ablation Studies}\label{sec:ablation}
%% ========================================================================

A sweep of $253$ runs over $66$ configurations on MNIST isolates subspace mode, the two
radii and the particle dimension; the temperature schedule, blockwise strategy,
schedule fragility and sensitivity plots are in Appendix~\ref{app:ablation}.

\paragraph{Ablation protocol.}
This sweep scores by maximum test accuracy over epochs rather than by the protocol of
Section~\ref{sec:setup}. The radius and dimension comparisons use a fixed epoch budget. Most
settings have five seeds; the heatmaps mark cells with only one seed. Maximum-test selection can still bias settings differently,
so even the ordering should be interpreted cautiously. These ablations are
not compared numerically with the primary validation-selected results.
Each ablation changes one setting in the deployed configuration
(Appendix~\ref{app:turbo-mode}). The controlled weighting study in
Section~\ref{sec:controlled-comparison} disables acceleration in a separate
fixed-subspace setup; differences between the studies cannot be attributed
to any single acceleration feature.

\paragraph{What the update rule contributes.}
With the same cost evaluator, entropic OT and row-wise softmax reach
$96.4\%$, while min-cost greedy and top-$k$ averaging reach about $12\%$
under this ablation's settings (Table~\ref{tab:ablation-summary}).
The plateau theorem alone does not explain these accuracy differences: it gives
zero displacement for all these deterministic equivariant rules on constant
rows. Column coupling has no resolved accuracy effect in this ablation
(Section~\ref{sec:ot-binding}). Reusing the transport plan across steps
is the most effective acceleration feature and the most fragile, since on a jumping
objective the transport matrix changes discontinuously between steps and a reused plan
answers a different problem; the five non-differentiable architecture benchmarks
therefore solve afresh at every step, as does theory mode everywhere.

\paragraph{Subspace and radius choices.}
Per-layer projection at rank $4$ has the highest mean among the four tested
subspace modes (Table~\ref{tab:ablation-summary}). The global projection
also has much larger variation across seeds. These comparisons do not
isolate the geometric cause of the differences. The radius sweeps in
Figure~\ref{fig:ablation-plate} favor larger configured multipliers over
the tested ranges. Their maximum-test metric does not establish an optimal
radius for validation-selected performance.

\begin{table}[t]
\centering
\caption{Ablation summary (MNIST, maximum test accuracy over epochs under the protocol above, so comparable within the table only). Subspace modes are read at fixed steps; particle dimension is read in the sweep's orthoplex configurations, and the simplex configurations every reported run uses return identical per-seed accuracies at both dimensions swept.}
\label{tab:ablation-summary}
\small
\begin{tabular}{@{}llc@{}}
\toprule
Choice & Setting & Test acc.\ (\%) \\
\midrule
\multirow{4}{*}{Subspace mode}
 & per-layer, rank 4 & $\mathbf{93.7 \pm 0.2}$ \\
 & single global projection & $86.1 \pm 17.2$ \\
 & full space & $84.2 \pm 0.2$ \\
 & adaptive & $71.5 \pm 1.3$ \\
\midrule
\multirow{3}{*}{Particle dimension (orthoplex)}
 & $d_p = 2$ & $\mathbf{95.7}$ \\
 & $d_p = 4$ & $94.9$ \\
 & $d_p = 8$ & $93.7$ \\
\midrule
\multirow{4}{*}{Update rule}
 & entropic OT & $\mathbf{96.4 \pm 0.2}$ \\
 & per-row softmax & $\mathbf{96.4 \pm 0.2}$ \\
 & min-cost greedy & ${\sim}12$ \\
 & top-$k$ averaging & ${\sim}12$ \\
\bottomrule
\end{tabular}
\end{table}

\paragraph{Particle dimension and schedules.}
The orthoplex ablation favors smaller $d_p$ at fixed steps. This comparison
changes both the number of particles and the number of vertices per particle.
It does not establish the same ordering for the deployed simplex. At $d_p=8$,
the simplex uses $9$ vertices instead of $16$, reducing evaluations per
fresh step by a factor of $1.78$; both frames satisfy the geometric
condition in Theorem~\ref{thm:piecewise-smooth}.

The tested cosine schedules reduce late-run accuracy by $2.24$ pp at
$102$K parameters and $2.46$ pp at $1$M parameters, with up to $4.5$ times
the reported standard deviation. These are observations under the ablation
protocol. The displacement envelope does not predict an accuracy difference.

\subsubsection{Binding of the Transport Constraint}\label{sec:ot-binding}

In the tested subspace diagnostic, the column-marginal residual is about
$3\times10^{-7}$ across the reported penalty sweep. The two solvers also
give similar accuracy across the rank sweep. Small residuals and similar
accuracy do not imply identical plans or barycentric updates. Full-space
MAX-SAT and RL are separate configurations.

Whether the constraint helps once it does bind is not settled by our data.
Appendix~\ref{app:ot-highparticle} gives both cases: a full-space MNIST run where
entropic OT leads softmax by $23$~pp, scored by maximum test accuracy over epochs and at
seed spreads of $\pm 11.0$ and $\pm 13.9$; and full-space
MAX-SAT at four instance sizes under the reported protocol, where the two differ by
between $-0.084$ and $+0.028$~pp with alternating sign. What the constraint supplies is
therefore the separation of Theorem~\ref{prop:plateau-freeze} and
Proposition~\ref{prop:ghost-drift}, which is proved, and not an accuracy gain, which is
not.

\subsection{Nonlinear Jump Experiments}\label{sec:jump-experiments}

Figure~\ref{fig:nonlinear-jumps} tests the finite-temperature predictions of
Section~\ref{sec:nonlinear-jump} on explicit losses. For planar jumps, we draw
4,096 Haar orthoplexes in each dimension $d_p\in\{2,3,8,16\}$ and reuse each
frame across 17 logarithmically spaced values of $J/\varepsilon$ from
$10^{-2}$ to $10^2$. Every realized drift agrees with
Eq.~\ref{eq:jump-softmax-drift} to within $4.5\times10^{-16}$ in float64.
The drift remains negative after the cost spread exceeds the temperature,
where the conditional saturation assumption no longer covers this example.

For curved jumps, let $\psi(x,y)=x+\kappa y^2/2$,
$g_-(x,y)=0.02y$, and $g_+(x,y)=2+0.03x+0.06y$, assigning the boundary
to $g_+$. Both branches are $G$-Lipschitz with $G=0.08$.
Starting at the origin, we use $\varepsilon=1$, $\ell=0.1$,
$\kappa\in\{0,0.5,2,8,20\}$, $\rho\in\{0.01,0.05,0.1,0.2\}$,
and independent uniform cost errors in $[-b,b]$ for
$b\in\{0,0.01,0.1,0.5\}$. Each of the 80 settings receives 4,096 fresh
rotations and noise draws from NumPy's generator with seed 20260913,
following the planar draws. All 173,274 trials satisfying the stated
sufficient conditions cross to the lower branch and decrease the loss.
Overall, 327,343 of 327,680 trials decrease it. The sufficient conditions
therefore exclude many successful steps; their failure does not predict an
increase. These simulations check the local mechanism and do not establish
the assumptions or performance of a trained network.

\begin{figure}[t]
\centering
\includegraphics[width=\linewidth]{nonlinear_jump_experiment}
\caption{Nonlinear softmax motion on explicit jumps, 4,096 frames per setting.
(a)~Mean inward normal drift with one sample standard deviation; the dotted
line marks $J=\varepsilon$. (b,c)~Curved-jump results at $\rho=0.2$,
$J=2$, $\varepsilon=1$, and $\ell=0.1$, colored by cost-error bound $b$.
(b)~Fraction meeting the sufficient conditions. (c)~Observed loss decrease
(solid, pointwise 95\% Wilson bands) and lower-side crossing (dashed);
the vertical axis is enlarged. These bands describe Monte Carlo sampling.}
\label{fig:nonlinear-jumps}
\end{figure}

\subsection{Controlled Weighting Comparison}\label{sec:controlled-comparison}

We compared softmax, linear, exponential-rank, and greedy weighting on the
hard-LIF MNIST model using identical candidate laws and independently tuned
step scales. All four arms completed validation tuning and ten final seeds.
Table~\ref{tab:controlled-weighting} summarizes all ten final seeds in each arm;
no run had a numerical failure.
The study used natural displacements, so weighting and realized step length
both contribute to the measured differences. It does not isolate a benefit
of softmax at a prescribed displacement length.

\begin{table}[t]
\centering
\caption{Controlled hard-LIF MNIST comparison, ten paired final seeds with
$10^7$ candidate evaluations per run. Accuracy is mean $\pm$ sample standard
deviation at the validation-selected checkpoint. Differences are softmax
minus the named rule, in percentage points, with 95\% paired percentile-bootstrap
intervals. Two-sided sign-flip $p$-values are unadjusted.}
\label{tab:controlled-weighting}
\small
\begin{tabular}{lccc}
\toprule
Weighting & Test accuracy (\%) & Difference [95\% interval] & $p$ \\
\midrule
Softmax & $74.49\pm9.73$ & $\mathrm{n/a}$ & $\mathrm{n/a}$ \\
Linear & $71.96\pm10.94$ & $2.53\ [0.20,4.83]$ & $0.082$ \\
Exponential rank & $78.42\pm10.10$ & $-3.93\ [-7.62,-0.18]$ & $0.088$ \\
Greedy & $75.89\pm9.27$ & $-1.40\ [-5.82,2.62]$ & $0.568$ \\
\bottomrule
\end{tabular}
\end{table}

Rank weighting had the highest mean accuracy. No available sign-flip test
rejected at the 5\% level. The bootstrap intervals and sign-flip tests use
different constructions; the intervals are not inversions of these tests.
The results do not establish general softmax superiority or equivalence
between the rules. This control evaluates the weighting rules within one fixed search representation
and candidate law; its outcome does not decompose the primary implementation's
acceleration and representation choices.

\paragraph{Data and randomization.}
We used the architecture, spike-count loss, and MNIST normalization from
Section~\ref{sec:setup}. A CPU PyTorch generator with seed 0 partitioned the
60K official training images into a 54K pool and 6K validation images.
A second permutation from that generator reserved a separate 1,024-image
holdout, leaving 52,976 update examples. The holdout was excluded from
updates and configuration selection. The official
10K test images were used only after validation selection. The split was
fixed across all seeds and methods. Minibatches contained 128 examples,
shuffled without replacement each epoch, with the final short batch retained.
Tuning used seeds 2027--2029; final runs used seeds 1000--1009. Initializations,
projection draws, minibatch order, and probe draws were paired across rules
within each seed, using separate random streams for data, projections, and probes.

\paragraph{Candidates and updates.}
We used a fixed per-layer rank-4 HybridSubspace with dense QR bases,
4,336 active coordinates, and two frozen coordinates in the incomplete
final block. Each of the 542 complete blocks had $d_p=8$ coordinates and
$V=16$ unit directions from a fresh Haar-rotated orthoplex, with $K=1$.
All rules used the same directions and minibatch at a given update index.
For costs $C_v$, directions $u_v$, and positive scale $\tau$, the rules were
\[
D_{\mathrm{softmax}}=\sum_v\mathrm{softmax}(-C/\tau)_v u_v,
\qquad
D_{\mathrm{linear}}=-\frac{1}{V\tau}\sum_v(C_v-\bar C)u_v,
\qquad
D_{\mathrm{rank}}=\sum_v\mathrm{softmax}(-r/\tau)_v u_v,
\]
where $r_v\in[0,1]$ is the average tied rank divided by $V-1$.
Greedy weighting averaged all minimum-cost vertices. Every rule returned
zero for constant costs. Linear weighting used signed coefficients rather
than a convex barycenter. Each block update was $\ell D$ with probes at
radius $\rho$. For softmax, these scales correspond to
$\varepsilon=\tau$, $r_p=2\rho/\tau$, and $r_s=\ell/\tau$.
Plan reuse, momentum, radius adaptation, projection refresh, and jitter were disabled.

\paragraph{Tuning, budgets, and selection.}
Softmax and linear weighting each used 27 configurations combining
$\rho\in\{10^{-3},10^{-2},10^{-1}\}$,
$\ell/\rho\in\{0.1,1,10\}$, and
$\tau\in\{10^{-3},10^{-2},10^{-1}\}$.
Rank weighting used the same radii and step ratios with
$\tau\in\{0.05,0.2,1\}$. Greedy used the same three radii and nine
step ratios $10^{k/4}$, $k=-4,\ldots,4$.
Each configuration received $2\times10^6$ candidate evaluations per tuning
seed; each final run received $10^7$. One update required 8,672 evaluations,
so tuning and final runs completed 230 and 1,153 updates, respectively,
leaving 5,440 and 1,184 evaluations unused. An update was admitted only if
its full candidate set fit the remaining budget.
All four arms therefore used the same tuning and final evaluation allowances.

Validation used all 6K images at initialization, after an update crossed
each additional $10^5$ candidate evaluations, and at the final update.
The selected checkpoint maximized validation accuracy, breaking ties by
lower validation loss and then earlier checkpoint. Configuration selection
maximized mean best validation accuracy across the three tuning seeds,
with ties broken by mean validation loss and then the configuration identifier.
Each final seed used its arm's selected configuration and evaluated the
selected checkpoint on test once. Validation selected
$(\rho,\ell,\tau)=(0.1,0.01,0.01)$ for softmax and linear weighting,
$(0.1,0.1,1)$ for rank weighting, and $(\rho,\ell)=(0.1,0.0177827941)$
for greedy. Final-seed test scores did not affect configuration selection.

\paragraph{Execution and uncertainty.}
Each run used a fresh process on the RTX 5090 with PyTorch 2.11.0 and CUDA 12.8,
float32, disabled TF32,
and the checkout's compiled forward evaluator with candidate chunks of 512.
No completed final run fell back to eager evaluation. Strict deterministic
algorithms were not enabled. Records include source hashes, software and
device information, split indices, configuration, seed, and validation
trajectories against updates, candidate evaluations, processed examples,
and elapsed time. These records do not establish a time-matched advantage.

All ten final seeds were retained. The numerical-failure rule specified the
last finite eligible validation checkpoint, or initialization if none existed;
no final run required this rule. Paired percentile-bootstrap
intervals used 10,000 resamples with NumPy seed 0. Tests enumerated all
$2^{10}$ sign assignments and assumed symmetry of paired differences under
the null. These three comparisons belong to a broader prespecified family
of 16 contrasts across five weighting rules, two candidate laws, and two
step normalizations. We report unadjusted tests and pointwise intervals
for the observed comparisons; no family-wide significance claim is made.
